\documentclass{scrbook}
\usepackage[utf8]{inputenc}

%General packages for basic document setup
\usepackage{amsmath,amsfonts, amsthm, graphicx}
\usepackage[margin=1.5cm,bottom=3cm, top=3cm]{geometry}

% Random text generator package
\usepackage{lipsum}

% Hyperlink package
\usepackage{hyperref}
\hypersetup{
  colorlinks=true,    %active colors for links
  linkcolor=blue,     %internal link color
  urlcolor=blue,      %external link color
  pdftitle={Template},%title of document (usually for pdf readers to display)
}

% Verbatism package : used for code display 
\usepackage{fancyvrb}
\usepackage{minted} %Interpret code and display colors aswell.

% Header and footer package
\usepackage{fancyhdr, lastpage}
% Header and footer parameters
% Some basic elements are \thepage for the number of the current page,
% \pageref{LastPage} for the number of the last page (add * to remove 
% the hyperlink.
% \rightmark for the current section, \leftmark for the current chapter,
\pagestyle{fancy}


%Chapter-Style : Only use with documentclass that require chaters !
\usepackage[Glenn]{fncychap}
\ChNameVar{\bfseries\Large\sffamily}
\ChTitleVar{\bfseries\Large\rmfamily}

%Colors package
\usepackage{xcolor}


% package for better refs
\usepackage[nameinlink]{cleveref}
% Some parameters to recognize tcbtheorem boxes as theorems, propositions, ...
\crefformat{theoremenv}{#2Théorème~#1#3}
\Crefformat{theoremenv}{#2Théorème~#1#3}

\crefformat{propositionenv}{#2Proposition~#1#3}
\Crefformat{propositionenv}{#2Proposition~#1#3}

\crefformat{lemmaenv}{#2Lemme~#1#3}
\Crefformat{lemmaenv}{#2Lemme~#1#3}

\crefformat{corollaryenv}{#2Corollaire~#1#3}
\Crefformat{corollaryenv}{#2Corollaire~#1#3}

\crefformat{definitionenv}{#2Définition~#1#3}
\Crefformat{definitionenv}{#2Définition~#1#3}

\crefformat{exampleenv}{#2Exemple~#1#3}
\Crefformat{exampleenv}{#2Exemple~#1#3}

\crefformat{remarkenv}{#2Remarque~#1#3}
\Crefformat{remarkenv}{#2Remarque~#1#3}

\newcommand{\creflastconjunction}{ et }

%Figure package
\usepackage{tikz}

%Boxes packages, for theorems, definitions, ...
%REQUIRES tikz in order to work !
\usepackage[most]{tcolorbox}
\tcbuselibrary{breakable}
% if only tcbset is used, it is not possible to make nested environments
% Hence the use of tcbsetforeverylayer
\tcbsetforeverylayer{enhanced jigsaw,breakable,fonttitle=\bfseries}

% ----------- New environments ----------

\usepackage{aliascnt}
% Theorems

\newtheoremstyle{theoremstyle}% name of the style to be used
  {\topsep}% measure of space to leave above the theorem. E.g.: 3pt
  {\topsep}% measure of space to leave below the theorem. E.g.: 3pt
  {\itshape}% name of font to use in the body of the theorem
  {0pt}% measure of space to indent
  {\bfseries}% name of head font
  {}% punctuation between head and body
  { }% space after theorem head; " " = normal interword space
  {
    \textcolor{red}
    {\thmname{#1} \thmnumber{#2}{\thmnote{#3}}}
  }

\theoremstyle{theoremstyle}
\newtheorem{theoremenv}{Théorème}[chapter]
\newenvironment{theorem}[1]{
  \begin{tcolorbox}[frame hidden, boxrule=0pt ,
    borderline west={2pt}{0pt}{red}, colback=red!10!white]{}
    \begin{theoremenv}{#1}
  }{
  \end{theoremenv}
\end{tcolorbox}
}

\newtheoremstyle{proofstyle}% name of the style to be used
  {\topsep}% measure of space to leave above the theorem. E.g.: 3pt
  {\topsep}% measure of space to leave below the theorem. E.g.: 3pt
  {}% name of font to use in the body of the theorem
  {0pt}% measure of space to indent
  {\bfseries}% name of head font
  {\textcolor{red!70!white}:}% punctuation between head and body
  { }% space after theorem head; " " = normal interword space
  {
    \textcolor{red!70!white}
    {\thmname{#1}}
  }

\theoremstyle{proofstyle}
\newtheorem{proofenv}{Preuve}
\renewenvironment{proof}{
  \vspace*{-0.74cm}
  \begin{tcolorbox}[frame hidden, boxrule=0pt ,
    borderline west={2pt}{0pt}{purple!60!white}, colback=red!10!white]{}
  \begin{proofenv}
  }{
  \qed
  \end{proofenv}
  \end{tcolorbox}
}

%Propositions
\newaliascnt{propositionenv}{theoremenv}
\newtheoremstyle{propositionstyle}% name of the style to be used
  {\topsep}% measure of space to leave above the proposition. E.g.: 3pt
  {\topsep}% measure of space to leave below the proposition. E.g.: 3pt
  {\itshape}% name of font to use in the body of the proposition
  {0pt}% measure of space to indent
  {\bfseries}% name of head font
  {}% punctuation between head and body
  { }% space after proposition head; " " = normal interword space
  {
    \textcolor{orange}
    {\thmname{#1} \thmnumber{#2}{\thmnote{#3}}}
  }

\theoremstyle{propositionstyle}
\newtheorem{propositionenv}[propositionenv]{Proposition}
\newenvironment{proposition}[1]{
  \begin{tcolorbox}[frame hidden, boxrule=0pt ,
    borderline west={2pt}{0pt}{orange}, colback=orange!10!white]{}
    \begin{propositionenv}{#1}
  }{
  \end{propositionenv}
\end{tcolorbox}
}
\aliascntresetthe{propositionenv}

%Lemmas
\newaliascnt{lemmaenv}{theoremenv}
\newtheoremstyle{lemmastyle}% name of the style to be used
  {\topsep}% measure of space to leave above the lemma. E.g.: 3pt
  {\topsep}% measure of space to leave below the lemma. E.g.: 3pt
  {\itshape}% name of font to use in the body of the lemma
  {0pt}% measure of space to indent
  {\bfseries}% name of head font
  {}% punctuation between head and body
  { }% space after lemma head; " " = normal interword space
  {
    \textcolor{yellow!60!black}
    {\thmname{#1} \thmnumber{#2}{\thmnote{#3}}}
  }

\theoremstyle{lemmastyle}
\newtheorem{lemmaenv}[lemmaenv]{Lemme}
\newenvironment{lemma}[1]{
  \begin{tcolorbox}[frame hidden, boxrule=0pt ,
    borderline west={2pt}{0pt}{yellow}, colback=yellow!10!white]{}
    \begin{lemmaenv}{#1}
  }{
  \end{lemmaenv}
\end{tcolorbox}
}
\aliascntresetthe{lemmaenv}

%Corollaries
\newaliascnt{corollaryenv}{theoremenv}
\newtheoremstyle{corollarystyle}% name of the style to be used
  {\topsep}% measure of space to leave above the corollary. E.g.: 3pt
  {\topsep}% measure of space to leave below the corollary. E.g.: 3pt
  {\itshape}% name of font to use in the body of the corollary
  {0pt}% measure of space to indent
  {\bfseries}% name of head font
  {}% punctuation between head and body
  { }% space after corollary head; " " = normal interword space
  {
    \textcolor{yellow!60!black}
    {\thmname{#1} \thmnumber{#2}{\thmnote{#3}}}
  }

\theoremstyle{corollarystyle}
\newtheorem{corollaryenv}[corollaryenv]{Corollaire}
\newenvironment{corollary}[1]{
  \begin{tcolorbox}[frame hidden, boxrule=0pt ,
    borderline west={2pt}{0pt}{yellow}, colback=yellow!10!white]{}
    \begin{corollaryenv}{#1}
  }{
  \end{corollaryenv}
\end{tcolorbox}
}
\aliascntresetthe{corollaryenv}

%Definitions
\newaliascnt{definitionenv}{theoremenv}
\newtheoremstyle{definitionstyle}% name of the style to be used
  {\topsep}% measure of space to leave above the definition. E.g.: 3pt
  {\topsep}% measure of space to leave below the definition. E.g.: 3pt
  {}% name of font to use in the body of the definition
  {0pt}% measure of space to indent
  {\bfseries}% name of head font
  {}% punctuation between head and body
  { }% space after definition head; " " = normal interword space
  {
    \textcolor{green!80!black}
    {\thmname{#1} \thmnumber{#2}{\thmnote{#3}}}
  }

\theoremstyle{definitionstyle}
\newtheorem{definitionenv}[definitionenv]{Définition}
\newenvironment{definition}[1]{
  \begin{tcolorbox}[frame hidden, boxrule=0pt ,
    borderline west={2pt}{0pt}{green}, colback=green!10!white]{}
    \begin{definitionenv}{#1}
  }{
  \end{definitionenv}
\end{tcolorbox}
}
\aliascntresetthe{definitionenv}

%Examples
\newaliascnt{exampleenv}{theoremenv}
\newtheoremstyle{examplestyle}% name of the style to be used
  {\topsep}% measure of space to leave above the example. E.g.: 3pt
  {\topsep}% measure of space to leave below the example. E.g.: 3pt
  {}% name of font to use in the body of the example
  {0pt}% measure of space to indent
  {\bfseries}% name of head font
  {}% punctuation between head and body
  { }% space after example head; " " = normal interword space
  {
    \textcolor{blue!80!white}
    {\thmname{#1} \thmnumber{#2}{\thmnote{#3}}}
  }

\theoremstyle{examplestyle}
\newtheorem{exampleenv}[exampleenv]{Exemple}
\newenvironment{example}[1]{
  \begin{tcolorbox}[frame hidden, boxrule=0pt ,
    borderline west={2pt}{0pt}{blue}, colback=blue!10!white]{}
    \begin{exampleenv}{#1}
  }{
  \end{exampleenv}
\end{tcolorbox}
}
\aliascntresetthe{exampleenv}

%Remarks
\newaliascnt{remarkenv}{theoremenv}
\newtheoremstyle{remarkstyle}% name of the style to be used
  {\topsep}% measure of space to leave above the remark. E.g.: 3pt
  {\topsep}% measure of space to leave below the remark. E.g.: 3pt
  {}% name of font to use in the body of the remark
  {0pt}% measure of space to indent
  {\bfseries}% name of head font
  {}% punctuation between head and body
  { }% space after remark head; " " = normal interword space
  {
    \textcolor{violet!80!white}
    {\thmname{#1} \thmnumber{#2}{\thmnote{#3}}}
  }

\theoremstyle{remarkstyle}
\newtheorem{remarkenv}[remarkenv]{Remarque}
\newenvironment{remark}{
  \begin{tcolorbox}[frame hidden, boxrule=0pt ,
    borderline west={2pt}{0pt}{violet}, colback=violet!10!white]{}
    \begin{remarkenv}{}
  }{
  \end{remarkenv}
\end{tcolorbox}
}
\aliascntresetthe{remarkenv}

% --- Exercises and solutions environments ---- 
%Use this package for correct use of exercise and answer environment.
\usepackage{exercise}

% Config of the exercise package
\renewcounter{Exercise}[chapter] %counter for the exercises that resets on each chapter.
\renewcounter{Answer}[chapter]

% Header for the exercises
\def\ExerciseHeader{\noindent\ExerciseHeaderDifficulty\textcolor{brown}{\textbf{Exercice \ExerciseHeaderNB} : \ExerciseHeaderTitle} \\ }

% Header for the answers
\def\AnswerHeader{\noindent\textcolor{brown}{\textbf{Solution de l'exercice \ExerciseHeaderNB}}\\ }

%Exercise environment
\newenvironment{exo}[1]{
  \stepcounter{Exercise}
  \begin{tcolorbox}[frame hidden, boxrule=0pt ,
    borderline west={2pt}{0pt}{brown}, colback=brown!10!white]{}
  \begin{Exercise}[title={#1},label={\the\value{chapter}-\the\value{Exercise}}, number={\the\value{Exercise}},
    counter={Exercise}]
  }{
  % Next line is to create a hyperref to the solution
  % [Solution \refAnswer{\ExerciseLabel}]
  \end{Exercise} 
\end{tcolorbox}
}
%Solution environment
\newenvironment{ans}{
  \begin{tcolorbox}[frame hidden, boxrule=0pt ,
    borderline west={2pt}{0pt}{brown!50!white}, colback=brown!5!white]{}
  \begin{Answer}[ref=\ExerciseLabel]
}{\end{Answer} 
\end{tcolorbox}
}

%%%%% USAGE ADVICE : ALWAYS USE SOLUTION ENVIRONMENT RIGHT AFTER AN EXERCISE ENVIRONMENT OTHERWISE,
%%%%% IT WILL BREAK THE NUMBERING SYSTEM.

% Line spacing package
\usepackage{setspace}
\onehalfspacing

% package used for multicolumns 
\usepackage{multicol}


\sloppy
\usepackage[final, nopatch=footnote]{microtype} % the option is to avoid the 
% Warning : Unable to apply patch 'footnote' 

% ----------- New commands ------------ %


\newcommand*{\blank}{\hspace*{1pt}}
\newcommand*{\MnR}[0]{\mathscr{M}_{n}(\mathbb{R})}
\newcommand*{\GLnR}[0]{\mathbf{GL}_{n}(\mathbb{R})}
\newcommand*{\transpose}[1]{{}^t #1}
\newcommand*{\GL}[0]{\textup{GL}}
\newcommand{\Alt}{\mathfrak{A}}
\newcommand*{\Sym}[0]{\mathfrak{S}}
\newcommand{\W}[0]{\textup{W}}
\newcommand{\Stab}[0]{\textup{Stab}}
\newcommand{\Fix}[0]{\textup{Fix}}
\newcommand*{\rg}[0]{\textup{rg}}
\newcommand*{\Im}[0]{\textup{Im}}
\newcommand*{\Ker}[0]{\textup{Ker}}
\newcommand*{\Coker}[0]{\textup{Coker}}
\newcommand*{\Tr}[0]{\textup{Tr}}
\newcommand*{\sinc}[0]{\textup{sinc}}
\newcommand*{\abs}[1]{\left| #1 \right|}
\newcommand*{\norm}[1]{\left\lVert #1 \right\rVert}
\newcommand*{\indicatrice}[0]{\mathds{1}}
\newcommand*{\Vect}[1]{\textup{Vect}\left\{#1\right\}}
